\hypertarget{objetivos}{%
\section{Objetivos}\label{objetivos}}

\begin{frame}{Objetivos}
\protect\hypertarget{objetivos-1}{}
\begin{itemize}
\tightlist
\item
  Motivar o uso de representa\c{c}\~{o}es intermediárias no pipeline de
  compila\~{c}\~{a}o e apresentar suas propriedades desejáveis.
\end{itemize}
\end{frame}

\begin{frame}{Objetivos}
\protect\hypertarget{objetivos-2}{}
\begin{itemize}
\tightlist
\item
  Definir a sintaxe e a sem\^{a}ntica operacional big-step da Representa\c{c}\~{a}o
  Intermediária em Árvore (IRT).
\end{itemize}
\end{frame}

\begin{frame}{Objetivos}
\protect\hypertarget{objetivos-3}{}
\begin{itemize}
\tightlist
\item
  Apresentar as regras de tradu\c{c}\~{a}o dirigida por sintaxe que convertem a
  AST de TWhile e TImp para a IRT.
\end{itemize}
\end{frame}

\hypertarget{introduuxe7uxe3o}{%
\section{Introdu\c{c}\~{a}o}\label{introduuxe7uxe3o}}

\begin{frame}{O Papel da Representa\c{c}\~{a}o Intermediária}
\protect\hypertarget{o-papel-da-representauxe7uxe3o-intermediuxe1ria}{}
\begin{itemize}
\tightlist
\item
  Converter a AST diretamente em código de máquina
\item
  Problemas:

  \begin{itemize}
  \tightlist
  \item
    Código de baixa qualidade
  \item
    Compilador fortemente acoplado \`{a} arquitetura
  \end{itemize}
\end{itemize}
\end{frame}

\begin{frame}{O Papel da Representa\c{c}\~{a}o Intermediária}
\protect\hypertarget{o-papel-da-representauxe7uxe3o-intermediuxe1ria-1}{}
\begin{itemize}
\tightlist
\item
  A \textbf{RI} desacopla \emph{front-end} de \emph{back-end}
\item
  Um \emph{front-end} pode gerar RI

  \begin{itemize}
  \tightlist
  \item
    A partir da RI pode-se suportar múltiplas arquiteturas
  \end{itemize}
\end{itemize}
\end{frame}

\begin{frame}{Propriedades Desejáveis de uma RI}
\protect\hypertarget{propriedades-desejuxe1veis-de-uma-ri}{}
\begin{itemize}
\tightlist
\item
  \textbf{Simplicidade}: \(\approx 15\) formas sintáticas (análise e
  transforma\c{c}\~{a}o mais fáceis)
\item
  \textbf{Independ\^{e}ncia de máquina}: sem detalhes de hardware
\end{itemize}
\end{frame}

\begin{frame}{Propriedades Desejáveis de uma RI}
\protect\hypertarget{propriedades-desejuxe1veis-de-uma-ri-1}{}
\begin{itemize}
\tightlist
\item
  \textbf{Independ\^{e}ncia de linguagem}: múltiplos \emph{front-ends}
  compartilham o mesmo \emph{back-end}.
\item
  \textbf{Suporte a transforma\c{c}\~{o}es}: facilitar análise de fluxo e
  otimiza\c{c}\~{o}es
\item
  \textbf{Proximidade do código-alvo}: opera\c{c}\~{o}es de memória e controle
  explícitos
\end{itemize}
\end{frame}

\hypertarget{a-irt-sintaxe}{%
\section{A IRT: Sintaxe}\label{a-irt-sintaxe}}

\begin{frame}{Express\~{o}es da IRT}
\protect\hypertarget{expressuxf5es-da-irt}{}
\[\begin{array}{lcll}
e & ::= & \mathbf{CONST}(n)  \\
  & \mid & \mathbf{TEMP}(t)  \\
  & \mid & \mathbf{NAME}(\ell) \\
  & \mid & \mathbf{BINOP}(\mathit{op}, e_1, e_2) \\
  & \mid & \mathbf{MEM}(e) \\
  & \mid & \mathbf{CALL}(e_f, e_1, \ldots, e_n) \\
  & \mid & \mathbf{ESEQ}(s, e)
\end{array}\]
\end{frame}

\begin{frame}{Express\~{o}es da IRT}
\protect\hypertarget{expressuxf5es-da-irt-1}{}
\begin{itemize}
\tightlist
\item
  Booleanos: \(\mathbf{CONST}(0)\) = falso, \(\mathbf{CONST}(1)\) =
  verdadeiro
\item
  \textbf{TEMP}: o alocador de registradores decide se vai para
  registrador ou pilha
\end{itemize}
\end{frame}

\begin{frame}{Comandos da IRT}
\protect\hypertarget{comandos-da-irt}{}
\[\begin{array}{lcll}
s & ::= & \mathbf{MOVE}(\mathbf{TEMP}(t), e) \\
  & \mid & \mathbf{MOVE}(\mathbf{MEM}(e_a), e) \\
  & \mid & \mathbf{EXP}(e) \\
  & \mid & \mathbf{SEQ}(s_1, s_2) \\
  & \mid & \mathbf{JUMP}(e) \\
  & \mid & \mathbf{CJUMP}(e, \ell_t, \ell_f) \\
  & \mid & \mathbf{LABEL}(\ell) \\
  & \mid & \mathbf{RETURN}(e_1, \ldots, e_n)
\end{array}\]
\end{frame}

\hypertarget{semuxe2ntica-operacional}{%
\section{Sem\^{a}ntica Operacional}\label{semuxe2ntica-operacional}}

\begin{frame}{Estado e Valores}
\protect\hypertarget{estado-e-valores}{}
O estado da IRT é um par \(\sigma = \langle \tau, \mu \rangle\):

\begin{itemize}
\tightlist
\item
  \(\tau : \mathit{Temp} \rightharpoonup \mathbb{Z}\)

  \begin{itemize}
  \tightlist
  \item
    Armazenamento de temporários
  \end{itemize}
\item
  \(\mu : \mathbb{Z} \rightharpoonup \mathbb{Z}\)

  \begin{itemize}
  \tightlist
  \item
    Memória (endere\c{c}os → valores)
  \end{itemize}
\end{itemize}

Todos os valores s\~{a}o inteiros de precis\~{a}o de máquina.
\end{frame}

\begin{frame}{Express\~{o}es}
\protect\hypertarget{expressuxf5es}{}
\textbf{Regras para express\~{o}es}: \(\sigma \vdash e \Downarrow v\):

\[
\begin{array}{c}
  \frac{}{\sigma \vdash \mathbf{CONST}(n) \Downarrow n} \\
  \frac{\sigma.\tau(t) = v}{\sigma \vdash \mathbf{TEMP}(t) \Downarrow v}\\
\end{array}
\]
\end{frame}

\begin{frame}{Express\~{o}es}
\protect\hypertarget{expressuxf5es-1}{}
\[
\begin{array}{c}
   \frac{\sigma \vdash e_1 \Downarrow v_1 \quad \sigma \vdash e_2 \Downarrow v_2}
        {\sigma \vdash \mathbf{BINOP}(\mathit{op}, e_1, e_2) \Downarrow \mathit{op}(v_1, v_2)}
\end{array}
\]
\end{frame}

\begin{frame}{Comandos}
\protect\hypertarget{comandos}{}
Resultado \(r\) de um comando:

\begin{itemize}
\tightlist
\item
  \(\langle \sigma', \blacksquare \rangle\): termina\c{c}\~{a}o normal
\item
  \(\langle \sigma', \uparrow (v_1,\ldots,v_k) \rangle\): retorno de
  fun\c{c}\~{a}o
\end{itemize}
\end{frame}

\begin{frame}{Comandos}
\protect\hypertarget{comandos-1}{}
\[
  \begin{array}{c}
  \frac{\sigma \vdash e \Downarrow v \quad \sigma' = \langle \tau[t \mapsto v], \mu \rangle}
       {P, \sigma \vdash \mathbf{MOVE}(\mathbf{TEMP}(t), e) \Downarrow \langle \sigma', \blacksquare \rangle}
  \end{array}
\]
\end{frame}

\begin{frame}{Comandos}
\protect\hypertarget{comandos-2}{}
\[
  \begin{array}{c}
    \frac{P, \sigma \vdash s_1 \Downarrow \langle \sigma', \blacksquare \rangle \quad P, \sigma' \vdash s_2 \Downarrow \langle \sigma'', r \rangle}{P, \sigma \vdash \mathbf{SEQ}(s_1, s_2) \Downarrow \langle \sigma'', r \rangle}
  \end{array}
\]

\begin{itemize}
\tightlist
\item
  Se \(s_1\) retorna, \(s_2\) \textbf{n\~{a}o é executado} (retorno
  propaga).
\end{itemize}
\end{frame}

\hypertarget{traduuxe7uxe3o-dirigida-por-sintaxe}{%
\section{Tradu\c{c}\~{a}o Dirigida por
Sintaxe}\label{traduuxe7uxe3o-dirigida-por-sintaxe}}

\begin{frame}{Express\~{o}es}
\protect\hypertarget{expressuxf5es-2}{}
\begin{itemize}
\tightlist
\item
  Literais e variáveis
\end{itemize}

\[\mathcal{E}\lbrack \mathtt{false} \rbrack = \mathbf{CONST}(0) \qquad \mathcal{E}\lbrack \mathtt{true} \rbrack = \mathbf{CONST}(1)\]

\[\mathcal{E}\lbrack n \rbrack = \mathbf{CONST}(n) \qquad \mathcal{E}\lbrack x \rbrack = \mathbf{TEMP}(x)\]
\end{frame}

\begin{frame}{Express\~{o}es}
\protect\hypertarget{expressuxf5es-3}{}
\begin{itemize}
\tightlist
\item
  Opera\c{c}\~{o}es:
\end{itemize}

\[\mathcal{E}\lbrack e_1 \oplus e_2 \rbrack = \mathbf{BINOP}(\widehat{\oplus},\; \mathcal{E}\lbrack e_1 \rbrack,\; \mathcal{E}\lbrack e_2 \rbrack)\]
\end{frame}

\begin{frame}[fragile]{Express\~{o}es}
\protect\hypertarget{expressuxf5es-4}{}
\begin{longtable}[]{@{}llll@{}}
\toprule\noalign{}
TWhile & IRT & TWhile & IRT \\
\midrule\noalign{}
\endhead
\passthrough{\lstinline!+!} & ADD & \passthrough{\lstinline!==!} & EQ \\
\passthrough{\lstinline!-!} & SUB & \passthrough{\lstinline!<!} & LT \\
\passthrough{\lstinline!*!} & MUL & \passthrough{\lstinline!\&\&!} &
AND \\
\bottomrule\noalign{}
\end{longtable}
\end{frame}

\begin{frame}{Comandos}
\protect\hypertarget{comandos-3}{}
\begin{itemize}
\tightlist
\item
  Declara\c{c}\~{a}o:
\end{itemize}

\[\begin{array}{c}
   \mathcal{S}\lbrack \mathbf{var}\ x : T = e \rbrack = \mathbf{MOVE}(\mathbf{TEMP}(x),\; \mathcal{E}\lbrack e \rbrack)
\end{array}\]
\end{frame}

\begin{frame}{Comandos}
\protect\hypertarget{comandos-4}{}
\begin{itemize}
\tightlist
\item
  Atribui\c{c}\~{a}o:
\end{itemize}

\[\begin{array}{c}
   \mathcal{S}\lbrack x := e \rbrack = \mathbf{MOVE}(\mathbf{TEMP}(x),\; \mathcal{E}\lbrack e \rbrack)
\end{array}\]
\end{frame}

\begin{frame}{Comandos}
\protect\hypertarget{comandos-5}{}
\begin{itemize}
\tightlist
\item
  Condicional:
\end{itemize}

\[\mathcal{S}\lbrack \mathbf{if}\ e\ \mathbf{then}\ b_1\ \mathbf{else}\ b_2 \rbrack =\]

\[\mathbf{SEQ}(\mathbf{CJUMP}(\mathcal{E}\lbrack e \rbrack, \ell_t, \ell_f),\; \mathbf{SEQ}(\mathbf{LABEL}(\ell_t),\; \ldots))\]
\end{frame}

\begin{frame}[fragile]{Comandos}
\protect\hypertarget{comandos-6}{}
\begin{itemize}
\tightlist
\item
  While:
\end{itemize}

\[\mathcal{S}\lbrack \mathbf{while}\ e\ \mathbf{do}\ B \rbrack =\]

\begin{lstlisting}
LABEL(loop_head)
CJUMP(E[[e]], loop_body, loop_exit)
LABEL(loop_body)
  S[[B]]
JUMP(NAME(loop_head))
LABEL(loop_exit)
\end{lstlisting}

\begin{itemize}
\tightlist
\item
  \(\ell_h\) = cabe\c{c}a; \(\ell_b\) = corpo; \(\ell_f\) = saída
\item
  A condi\c{c}\~{a}o é reavaliada a cada itera\c{c}\~{a}o
\end{itemize}
\end{frame}

\begin{frame}[fragile]{Implementa\c{c}\~{a}o: A M\^{o}nada CgM}
\protect\hypertarget{implementauxe7uxe3o-a-muxf4nada-cgm}{}
\begin{lstlisting}[language=Haskell]
data CgState = CgState
  { freshCounter :: Int
  , emitted      :: [Stmt]
  }

type CgM = StateT CgState (ExceptT String Identity)

freshLabel :: String -> CgM Label
freshLabel prefix = do
    n <- gets freshCounter
    modify (\s -> s { freshCounter = n + 1 })
    return (prefix ++ "_" ++ show n)

emit :: Stmt -> CgM ()
emit s = modify (\st -> st { emitted = emitted st ++ [s] })
\end{lstlisting}
\end{frame}

\begin{frame}[fragile]{Compilando While}
\protect\hypertarget{compilando-while}{}
\begin{lstlisting}[language=Haskell]
compileTWhile :: [Stmt] -> CgM ()
compileTWhile stmts = mapM_ compileStmt stmts

compileStmt (SWhile cond body) = do
    lHead <- freshLabel "loop_head"
    lBody <- freshLabel "loop_body"
    lExit <- freshLabel "loop_exit"
    emit (LABEL lHead)
    cExpr <- compileExp cond
    emit (CJUMP cExpr lBody lExit)
    emit (LABEL lBody)
    mapM_ compileStmt body
    emit (JUMP (NAME lHead))
    emit (LABEL lExit)
\end{lstlisting}
\end{frame}

\hypertarget{conclusuxe3o}{%
\section{Conclus\~{a}o}\label{conclusuxe3o}}

\begin{frame}{Conclus\~{a}o}
\protect\hypertarget{conclusuxe3o-1}{}
\begin{itemize}
\tightlist
\item
  A \textbf{IRT} separa express\~{o}es (produzem valor) de comandos
  (produzem efeitos)
\item
  \textbf{7 express\~{o}es}: CONST, TEMP, NAME, BINOP, MEM, CALL, ESEQ
\item
  \textbf{7 comandos}: MOVE, EXP, SEQ, JUMP, CJUMP, LABEL, RETURN
\end{itemize}
\end{frame}

\begin{frame}{Conclus\~{a}o}
\protect\hypertarget{conclusuxe3o-2}{}
\begin{itemize}
\tightlist
\item
  A sem\^{a}ntica big-step define o estado como par
  \(\langle \tau, \mu \rangle\)
\item
  A tradu\c{c}\~{a}o \textbf{dirigida por sintaxe} converte AST na IRT
  sistematicamente
\end{itemize}
\end{frame}
